%% =====================================================================
%%  metadata.tex
%%  Khai báo toàn bộ thông tin bài báo: tiêu đề, tác giả, đơn vị,
%%  tóm tắt, từ khóa, tác giả liên hệ, dòng "Trích dẫn bài báo này".
%%
%%  Các lệnh [vi] dùng cho trang bìa tiếng Việt (trang 1).
%%  Các lệnh [en] dùng cho trang bìa tiếng Anh (trang cuối).
%%  Nếu tắt bìa tiếng Anh (option noenglishcover) thì phần [en]
%%  vẫn có thể để nguyên, không ảnh hưởng gì.
%%
%%  LƯU Ý: file này chỉ khai báo, KHÔNG in gì ra trang.
%%         Nó được \input ngay sau \documentclass.
%% =====================================================================

%% ---------------------------------------------------------------------
%% 1. TIÊU ĐỀ
%%    - Viết hoa đúng chính tả, không cần IN HOA toàn bộ:
%%      class tự xử lý phần trình bày.
%%    - Công thức trong tiêu đề dùng $...$ như bình thường.
%% ---------------------------------------------------------------------
\jsttitlevi{%
  Lập kế hoạch đường đi lai ghép cho robot tự hành
  trong môi trường động sử dụng Voronoi có trọng số,
  \texorpdfstring{$\beta$}{β}-complexes và giám sát logic mờ%
}
\jsttitleen{%
  Hybrid Path Planning for Autonomous Mobile Robots
  Using Weighted Voronoi Diagrams,
  \texorpdfstring{$\beta$}{β}-complexes, and Fuzzy Logic Supervision%
}

%% ---------------------------------------------------------------------
%% 2. LOẠI BÀI (in ở góc phải header)
%% ---------------------------------------------------------------------
\jstarticletypevi{Bài nghiên cứu}
\jstarticletypeen{Research Article}

%% ---------------------------------------------------------------------
%% 3. TÁC GIẢ
%%    Cú pháp: \jstauthorvi[chỉ số đơn vị]{Họ và tên}
%%    - Chỉ số đơn vị in thành superscript, khớp với \jstaffiliationvi.
%%    - Bỏ [] nếu tác giả chỉ có một đơn vị.
%%    - Tác giả liên hệ (có dấu *) phải khai báo thêm ở mục 5.
%% ---------------------------------------------------------------------
\jstauthorvi[1,2]{Hoàng Minh Anh Tú}
\jstauthorvi[1]{Trần Văn Hoài}
\jstauthorvi[2,*]{Phan Thành An}

\jstauthoren[1,2]{Hoang Minh Anh Tu}
\jstauthoren[1]{Tran Van Hoai}
\jstauthoren[2,*]{Phan Thanh An}

%% ---------------------------------------------------------------------
%% 4. ĐƠN VỊ CÔNG TÁC (cột trái của bìa, chữ 7 pt nghiêng)
%%    Cú pháp: \jstaffiliationvi{chỉ số}{nội dung}
%%    Dùng \\ để xuống dòng trong cùng một đơn vị.
%%    Thứ tự khai báo = thứ tự hiển thị.
%% ---------------------------------------------------------------------
\jstaffiliationvi{1}{%
  Khoa Khoa học và Kỹ thuật Máy tính, Trường Đại học Bách khoa,\\
  Đại học Quốc gia Thành phố Hồ Chí Minh, Việt Nam}

\jstaffiliationvi{2}{%
  Khoa Khoa học và Kỹ thuật, Trường Đại học Quốc tế,\\
  Đại học Quốc gia Thành phố Hồ Chí Minh, Việt Nam}

\jstaffiliationen{1}{%
  Faculty of Computer Science and Engineering,
  Ho Chi Minh City University of Technology,\\
  Viet Nam National University Ho Chi Minh City, Vietnam}

\jstaffiliationen{2}{%
  Faculty of Engineering,
  International University,\\
  Viet Nam National University Ho Chi Minh City, Vietnam}

%% ---------------------------------------------------------------------
%% 5. TÁC GIẢ LIÊN HỆ (khối dưới đơn vị, có đường kẻ ngang phía trên)
%%    Cú pháp: \jstcorrespondingvi{Tên}{Địa chỉ}{Email}
%% ---------------------------------------------------------------------
\jstcorrespondingvi{Phan Thành An}%
  {Khoa Khoa học và Kỹ thuật, Trường Đại học Quốc tế, ĐHQG-HCM}%
  {ptan@hcmiu.edu.vn}

\jstcorrespondingen{Phan Thanh An}%
  {Faculty of Engineering, International University, VNU-HCM}%
  {ptan@hcmiu.edu.vn}

%% ---------------------------------------------------------------------
%% 6. TÓM TẮT + TỪ KHÓA
%%    Tóm tắt 9 pt, từ khóa 10 pt, nằm trong khung nền xanh nhạt.
%%    Viết thành một đoạn văn liền, không xuống dòng, không bullet.
%% ---------------------------------------------------------------------
\jstabstractvi{%
  Bài báo trình bày một khung lập kế hoạch đường đi lai ghép cho robot
  tự hành trong môi trường có vật cản động. Khung gồm ba lớp: biểu diễn
  hình học bằng Voronoi có trọng số và $\beta$-complexes, lấy mẫu
  đường đi bằng RRT* có hướng dẫn, và giám sát an toàn bằng logic mờ.
  Kết quả thử nghiệm trên robot vi sai cho thấy thời gian lập kế hoạch
  giảm đáng kể so với phương pháp cơ sở trong khi vẫn giữ được tính
  không va chạm của quỹ đạo.%
}
\jstkeywordsvi{%
  lập kế hoạch đường đi; Voronoi có trọng số; $\beta$-complexes;
  RRT*; logic mờ; robot tự hành%
}

\jstabstracten{%
  This paper presents a hybrid path-planning framework for autonomous
  robots operating in dynamic environments. The framework has three
  layers: geometric representation by weighted Voronoi diagrams and
  $\beta$-complexes, guided sampling by RRT*, and safety supervision
  by fuzzy logic. Experiments on a differential-drive robot show that
  planning time is substantially reduced compared with baseline
  planners while collision-free trajectories are preserved.%
}
\jstkeywordsen{%
  path planning; weighted Voronoi diagram; $\beta$-complexes;
  RRT*; fuzzy logic; autonomous robot%
}

%% ---------------------------------------------------------------------
%% 7. DÒNG "TRÍCH DẪN BÀI BÁO NÀY" (ô xanh ở đáy bìa)
%%    Điền theo đúng định dạng của tạp chí khi đã có số phát hành.
%% ---------------------------------------------------------------------
\jstcitationvi{%
  Tú HMA, Hoài TV, An PT. Lập kế hoạch đường đi lai ghép cho robot
  tự hành trong môi trường động sử dụng Voronoi có trọng số,
  $\beta$-complexes và giám sát logic mờ.
  Tạp chí Khoa học và Công nghệ ĐHQG-HCM, chuyên san Kỹ thuật và
  Công nghệ. 2026;Tập(So):trang.%
}
\jstcitationen{%
  Tu HMA, Hoai TV, An PT. Hybrid path planning for autonomous mobile
  robots using weighted Voronoi diagrams, $\beta$-complexes, and
  fuzzy logic supervision.
  VNUHCM Journal of Engineering and Technology. 2026;Vol(Issue):pages.%
}
